\documentclass[11pt,letterpaper]{article}

% ------------------------------------------------------------------
% Paquetes
% ------------------------------------------------------------------
\usepackage[utf8]{inputenc}
\usepackage[T1]{fontenc}
\usepackage[spanish,es-noshorthands]{babel}
\usepackage{newtxtext,newtxmath}          % Times New Roman
\usepackage[margin=2.5cm]{geometry}
\usepackage{graphicx}
\usepackage{float}
\usepackage{caption}
\usepackage{enumitem}
\usepackage{listings}
\usepackage{titlesec}
\usepackage{fancyhdr}
\usepackage[hidelinks]{hyperref}

% ------------------------------------------------------------------
% Estilo
% ------------------------------------------------------------------
\titleformat{\section}{\large\bfseries}{}{0em}{}[\vspace{-0.6em}\rule{\linewidth}{0.4pt}]
\titleformat{\subsection}{\normalsize\bfseries}{}{0em}{}
\titlespacing*{\section}{0pt}{1.2em}{0.8em}
\titlespacing*{\subsection}{0pt}{1em}{0.4em}

\pagestyle{fancy}
\fancyhf{}
\fancyhead[L]{\small DOPO -- Laboratorio 6: Persistencia}
\fancyhead[R]{\small Espitia -- Salamanca}
\fancyfoot[C]{\small\thepage}
\renewcommand{\headrulewidth}{0.4pt}

\captionsetup{font=small,labelfont=bf}
\setlength{\parindent}{0pt}
\setlength{\parskip}{0.5em}
\setlist[enumerate]{leftmargin=*,itemsep=0.3em}
\setlist[itemize]{itemsep=0.1em}

% Enunciado y respuesta
\newcommand{\pregunta}[1]{\textit{#1}\par}
\newcommand{\resp}{\textbf{R/}~}

% Espacio para captura: si el archivo existe en imagenes/ se incluye,
% si no, se deja un recuadro con la ruta esperada.
\newcommand{\captura}[3][0.75]{%
  \begin{figure}[H]
    \centering
    \IfFileExists{imagenes/#2}{%
      \includegraphics[width=#1\linewidth]{imagenes/#2}%
    }{%
      \fbox{\parbox[c][4.5cm][c]{#1\linewidth}{\centering
        \small Captura pendiente\\[4pt]\texttt{imagenes/#2}}}%
    }
    \caption{#3}
  \end{figure}}

% Código
\lstset{
  basicstyle=\ttfamily\small,
  frame=single,
  keywordstyle=\bfseries,
  commentstyle=\itshape,
  breaklines=true, columns=fullflexible, showstringspaces=false,
  xleftmargin=4pt, xrightmargin=4pt, aboveskip=0.8em, belowskip=0.8em,
  literate={á}{{\'a}}1 {é}{{\'e}}1 {í}{{\'i}}1 {ó}{{\'o}}1 {ú}{{\'u}}1
           {ñ}{{\~n}}1 {Á}{{\'A}}1 {É}{{\'E}}1 {Í}{{\'I}}1 {Ó}{{\'O}}1 {Ú}{{\'U}}1
}
\lstdefinestyle{java}{language=Java}
\lstdefinestyle{txt}{language={}, numbers=left, numberstyle=\tiny, numbersep=6pt, xleftmargin=16pt}

% ==================================================================
\begin{document}

% ------------------------------------------------------------------
% Portada
% ------------------------------------------------------------------
\begin{titlepage}
  \centering
  \vspace*{2cm}
  {\large\textsc{Escuela Colombiana de Ingeniería Julio Garavito}\par}
  \vspace{0.4cm}
  {\normalsize Desarrollo Orientado por Objetos (DOPO) -- Grupo 501\par}
  \vspace{3cm}
  \rule{\linewidth}{0.8pt}\par
  \vspace{0.4cm}
  {\Huge\bfseries Laboratorio 6\par}
  \vspace{0.3cm}
  {\LARGE Persistencia\par}
  \vspace{0.4cm}
  \rule{\linewidth}{0.8pt}\par
  \vspace{3cm}
  {\large\textbf{Integrantes}\par}
  \vspace{0.2cm}
  {\large Juan David Espitia\par Cristian Camilo Salamanca\par}
  \vspace{1.2cm}
  {\large\textbf{Profesora}\par}
  \vspace{0.2cm}
  {\large María Irma Díaz Rozo\par}
  \vfill
  {\large Bogotá D.C., 2026-2\par}
\end{titlepage}

\tableofcontents
\newpage

% ==================================================================
\section{Parte I. Creando la maqueta}
% ==================================================================

\subsection{A. Safari inicial}

\begin{enumerate}
  \item \pregunta{¿Cuál ambiente seleccionaron? ¿Por qué?}
  \resp Inicialmente se usó BlueJ por un inconveniente con la importación de fuentes. Sin embargo, se eligió VS Code como ambiente principal, ya que permite un desarrollo orientado a objetos más profesional: navegación entre clases, refactorización y ejecución integrada de pruebas JUnit.

  \item \pregunta{Ejecuten el programa y capturen la pantalla.}
  \captura{I-A-2-ejecucion.png}{Ejecución inicial de EcoSafari.}
  \captura{I-A-2-ejecucion-2.png}{Ejecución de EcoSafari después de algunos clics.}
\end{enumerate}

\subsection{B. Maqueta: MVC}

\begin{enumerate}
  \item \textbf{Modelo.} En la fachada \texttt{EcoSafari} se crearon los métodos \texttt{open}, \texttt{saveAs}, \texttt{import\_} y \texttt{exportAs}, que reciben un \texttt{File} y propagan una \texttt{EcoSafariException}:
\begin{lstlisting}[style=java]
public void open(File file) throws EcoSafariException {
    String fileName = (file != null) ? file.getName() : "";
    throw new EcoSafariException("Opción open en construcción. Archivo " + fileName);
}
\end{lstlisting}
  \captura{I-B-1-modelo.png}{Métodos de entrada-salida en la fachada del dominio.}

  \item \textbf{Vista.} El método \texttt{prepareViewMenu} construye el menú \textit{Archivo} con las seis opciones y sus separadores.
  \captura{I-B-2-menu.png}{Método prepareViewMenu.}
  \captura[0.55]{I-B-2-menu-abierto.png}{Menú Archivo.}

  \item \textbf{Controlador.} El método \texttt{prepareControllerMenu} asocia cada opción con su método: \texttt{optionNew}, \texttt{optionOpen}, \texttt{optionSaveAs}, \texttt{optionImport}, \texttt{optionExportAs} y \texttt{optionExit}.
  \begin{enumerate}[label=\alph*)]
    \item \textbf{Salir:} \texttt{optionExit} termina la aplicación con \texttt{System.exit(0)}, sin confirmación.
    \item \textbf{Nuevo:} \texttt{optionNew} reemplaza el safari actual por uno nuevo y repinta la vista.
    \captura{I-B-3b-nuevo.png}{Resultado de la opción Nuevo.}
    \item \textbf{Entrada-salida:} cada opción abre un \texttt{JFileChooser}, llama al método del modelo y muestra la excepción en un \texttt{JOptionPane}.
    \captura[0.6]{I-B-3c-abrir.png}{Abrir: mensaje de opción en construcción.}
    \captura[0.6]{I-B-3c-salvar.png}{Salvar como: mensaje de opción en construcción.}
    \captura[0.6]{I-B-3c-importar.png}{Importar: mensaje de opción en construcción.}
    \captura[0.6]{I-B-3c-exportar.png}{Exportar como: mensaje de opción en construcción.}
  \end{enumerate}
\end{enumerate}

% ==================================================================
\section{Parte II. Implementando persistencia}
% ==================================================================

\subsection{A. Salvar y abrir}

\begin{enumerate}
  \item Las versiones de maqueta se conservaron como \texttt{open00} y \texttt{saveAs00}.

  \item \textbf{saveAs.} Serializa el safari completo con \texttt{ObjectOutputStream}. Para ello \texttt{EcoSafari} y todas las entidades implementan \texttt{Serializable} y declaran \texttt{serialVersionUID}.
\begin{lstlisting}[style=java]
try (ObjectOutputStream out = new ObjectOutputStream(new FileOutputStream(file))) {
    out.writeObject(this);
} catch (IOException e) {
    throw new EcoSafariException("el siguiente archivo no ha sido posible guardarlo " + file.getName());
}
\end{lstlisting}

  \item \pregunta{¿El archivo se creó en el disco? ¿Cuánto espacio ocupa?}
  \resp Sí. \texttt{safariOne.dat} ocupa aproximadamente 2.3~KB (2\,311 bytes). El tamaño varía ligeramente entre ejecuciones porque la simulación tiene un componente aleatorio (\texttt{Ground}).
  \captura{II-A-3-safariOne-dat.png}{Archivo safariOne.dat en el disco.}

  \item \textbf{open.} Lee el objeto con \texttt{ObjectInputStream} y reemplaza el tablero actual por el leído.

  \item \pregunta{Prueba de aceptación: iniciar, crear un nuevo estado y abrir safariOne.dat.}
  \captura{II-A-5-nuevo-estado.png}{Nuevo estado antes de abrir.}
  \captura{II-A-5-abierto.png}{Estado restaurado desde safariOne.dat.}
\end{enumerate}

\subsection{B. Exportar e importar}

\begin{enumerate}
  \item Las versiones de maqueta se conservaron como \texttt{import00} y \texttt{exportAs00}.

  \item \textbf{exportAs.} Recorre el tablero y escribe una línea por entidad con el formato \texttt{Tipo fila, columna}, usando \texttt{getClass().getSimpleName()}.

  \item \pregunta{Exportar como safariOne.txt, editarlo y analizar. ¿Qué pasó?}
  \resp Se generó un archivo de texto legible, con una línea por cada entidad. Las celdas donde una entidad estaba sobre el suelo (por ejemplo, \texttt{Grass} sobre \texttt{Ground}) solo aparecen con la entidad visible. El archivo guarda tipo y posición, pero no el estado interno (energía).
  \captura{II-B-3-safariOne-txt.png}{Contenido de safariOne.txt exportado.}

  \item \textbf{import\_.} Lee cada línea, la separa con \texttt{trim} y \texttt{split} y crea la entidad correspondiente en la posición indicada.

  \item \pregunta{Exportar a safariOther.txt, salir, entrar, Nuevo e importar. ¿Qué resultado obtuvieron?}
  \resp Se recuperó el mismo tablero: mismos tipos y posiciones. Las entidades inician con su estado por defecto (energía 100), ya que el archivo de texto no lo guarda.
  \captura{II-B-5-safariOther.png}{Safari después de importar safariOther.txt.}

  \item \pregunta{Escribir un safariOne.txt correcto e importarlo. ¿Qué resultado obtuvieron?}
  \resp El tablero se construyó exactamente con las entidades del archivo.
\begin{lstlisting}[style=txt]
Elephant 5, 5
Elephant 10, 10
Shrub 6, 6
Storm 0, 20
SavannaElephant 8, 8
Chameleon 2, 0
Grass 19, 2
Zebra 20, 2
Lion 22, 2
Ground 18, 0
...
\end{lstlisting}
  \captura{II-B-6-safariOne-importado.png}{Safari importado desde safariOne.txt.}

  \item Se eliminó el método \texttt{someEntities()}. El constructor crea un safari vacío y los elementos solo se agregan mediante archivos.
\end{enumerate}

\subsection{C. Comportamiento}

\begin{enumerate}
  \item \pregunta{Tres clics, guardar y abrir.}
  \resp Al abrir, el safari continúa exactamente donde quedó: posiciones, energía y estado de cada entidad.
  \captura{II-C-1-guardar-abrir.png}{Safari después de guardar y abrir.}

  \item \pregunta{Tres clics, exportar e importar.}
  \resp Se recuperan los tipos y posiciones, pero las entidades reinician su estado interno.
  \captura{II-C-2-exportar-importar.png}{Safari después de exportar e importar.}

  \item \pregunta{¿Qué diferencias ven entre 1 y 2?}
  \resp Guardar crea un archivo binario con el objeto completo, incluido el estado interno de cada entidad. Exportar crea un archivo de texto, editable a mano, que solo describe qué entidad hay y dónde. Por eso, al importar, cada entidad se construye de nuevo con sus valores iniciales, mientras que al abrir la simulación continúa sin cambios.
\end{enumerate}

% ==================================================================
\section{Parte III. Fortaleciendo la persistencia}
% ==================================================================

\subsection{A. Guardar y abrir}

\begin{enumerate}
  \item Las versiones con mensaje general se conservaron como \texttt{open01} y \texttt{saveAs01}.

  \item Antes de leer o escribir se valida el archivo, y cada caso tiene su propio mensaje:
  \begin{itemize}
    \item Archivo nulo: \textit{No se ha especificado ningún archivo para abrir/guardar.}
    \item Archivo inexistente (\texttt{open}): \textit{el siguiente archivo no existe.}
    \item Directorio: \textit{el siguiente archivo en realidad es un directorio.}
    \item Extensión distinta de \texttt{.dat}: \textit{el siguiente archivo no es un .dat.}
    \item Contenido inválido (\texttt{open}): \textit{no ha sido posible abrirlo por corrupción o incompatibilidad con EcoSafari.}
    \item Falla de escritura (\texttt{saveAs}): \textit{no ha sido posible guardarlo por error de escritura o permisos.}
  \end{itemize}

  \item \pregunta{Prueba de aceptación de un nuevo mensaje.}
  \resp Se intentó abrir un archivo inexistente. Resultado: \textit{``el siguiente archivo no existe: noexiste.dat''}.
  \captura[0.6]{III-A-3-mensaje.png}{Mensaje de error detallado al abrir.}
\end{enumerate}

\subsection{B. Exportar e importar}

\begin{enumerate}
  \item Las versiones con mensaje general se conservaron como \texttt{import01} y \texttt{exportAs01}.

  \item Se aplicaron las mismas validaciones de archivo (nulo, inexistente, directorio y extensión \texttt{.txt}). Además, \texttt{import\_} distingue líneas con datos incompletos, coordenadas no numéricas y tipos de entidad no reconocidos.

  \item \pregunta{Prueba de aceptación de un nuevo mensaje.}
  \resp Se intentó exportar a un archivo con extensión \texttt{.dat}. Resultado: \textit{``el siguiente archivo no es un .txt: a.dat''}.
  \captura[0.6]{III-B-3-mensaje.png}{Mensaje de error detallado al exportar.}
\end{enumerate}

\subsection{Pruebas de unidad}

Las pruebas de persistencia están en \texttt{EcoSafariPersistenceTest} (24 pruebas) y cubren guardar y abrir, exportar e importar, cada error del minicompilador y la diferencia de estado entre ambos mecanismos: \texttt{open} conserva la energía y \texttt{import\_} la reinicia. El proyecto completo suma 53 pruebas, todas exitosas.

\captura[0.6]{III-pruebas.png}{Ejecución de las pruebas de unidad.}

% ==================================================================
\section{Parte IV. Bono: perfeccionando importar}
% ==================================================================

\subsection{A. Minicompilador}

\begin{enumerate}
  \item La versión anterior se conservó como \texttt{import02}.

  \item Cada error indica la línea, la palabra que lo produce y la causa, con el formato:
  \begin{center}
    \texttt{Línea \textit{n}: error en '\textit{palabra}'. Causa: \textit{descripción}.}
  \end{center}
  Además, si alguna línea falla se conserva el tablero anterior, de modo que un archivo con errores no deja el safari a medias.

  \item \pregunta{safariErr.txt con mínimo tres errores diferentes.}
\begin{lstlisting}[style=txt]
Dragon 5, 5
Elephant 10, abc
Lion 50, 50
Bush 15, 15
Storm 5, 5
\end{lstlisting}
  \resp El archivo se fue corrigiendo línea a línea con ayuda del importador, que reportó en orden: la clase \texttt{Dragon} no existe; la columna \texttt{abc} no es un número entero; las coordenadas \texttt{50, 50} están fuera del tablero (0 a 24); y la clase \texttt{Bush} no existe, pues en la aplicación se llama \texttt{Shrub}.

  \captura[0.6]{IV-A-3-error-1.png}{Error 1: clase inexistente.}
  \captura[0.6]{IV-A-3-error-2.png}{Error 2: coordenada no numérica.}
  \captura[0.6]{IV-A-3-error-3.png}{Error 3: coordenadas fuera de límites.}
  \captura[0.6]{IV-A-3-error-4.png}{Error 4: nombre de clase incorrecto.}
\end{enumerate}

\subsection{B. Minicompilador flexible}

\begin{enumerate}
  \item La versión anterior se conservó como \texttt{import03}.

  \item El \texttt{switch} de tipos se reemplazó por reflexión: la clase se busca por su nombre y se invoca su constructor \texttt{(EcoSafari, int, int)}. Así cualquier entidad futura se importa sin modificar el método.
\begin{lstlisting}[style=java]
try {
    Class.forName("domain." + tipo)
         .getConstructor(EcoSafari.class, int.class, int.class)
         .newInstance(this, i, j);
} catch (ClassNotFoundException ce) {
    throw new EcoSafariException("Línea " + numeroLinea + ": error en '" + tipo
        + "'. Causa: la clase descrita en el archivo no existe.");
} catch (NoSuchMethodException | InstantiationException
         | IllegalAccessException | InvocationTargetException e) {
    throw new EcoSafariException("Línea " + numeroLinea + ": error en '" + tipo
        + "'. Causa: la clase descrita tiene inconvenientes en su constructor.");
}
\end{lstlisting}

  \item \pregunta{¿exportAs necesita cambiar para ser flexible?}
  \resp No. \texttt{exportAs} obtiene el tipo con \texttt{getClass().getSimpleName()}, por lo que ya funciona con cualquier clase nueva sin modificaciones.

  \item No aplica, ya que \texttt{exportAs} no requirió cambios.

  \item \pregunta{safariFlex.txt para probar la flexibilidad.}
\begin{lstlisting}[style=txt]
Elephant 5, 5
Lion 8, 8
Zebra 12, 12
Chameleon 3, 3
Shrub 15, 15
Storm 2, 20
Grass 19, 2
Ground 18, 0
SavannaElephant 10, 10
Dragon 1, 1
\end{lstlisting}
  \resp Las nueve primeras clases se reconocen sin estar escritas en el código del importador. La línea 10 produce el error significativo: \textit{``Línea 10: error en 'Dragon'. Causa: la clase descrita en el archivo no existe.''}
  \captura[0.6]{IV-B-5-safariFlex.png}{Error significativo al importar safariFlex.txt.}
\end{enumerate}

% ==================================================================
\section{Retrospectiva}
% ==================================================================

\begin{enumerate}
  \item \pregunta{¿Cuál fue el tiempo total invertido por cada uno? (Horas)}
  \resp Juan David Espitia: \textit{[completar]} horas. Cristian Camilo Salamanca: \textit{[completar]} horas.

  \item \pregunta{¿Cuál es el estado actual del laboratorio? ¿Por qué?}
  \resp Completo, incluido el bono. Todas las opciones del menú funcionan, los errores se reportan en detalle y las 53 pruebas de unidad pasan.

  \item \pregunta{¿Cuál práctica XP fue la más útil? ¿Por qué?}
  \resp La propiedad colectiva del código: ambos integrantes conocían y podían modificar cualquier clase, lo que facilitó corregir errores en el dominio y en la interfaz sin depender del otro.

  \item \pregunta{¿Cuál consideran fue el mayor logro? ¿Por qué?}
  \resp El minicompilador flexible: con reflexión, el importador reconoce cualquier entidad futura y reporta los errores con línea, palabra y causa.

  \item \pregunta{¿Cuál fue el mayor problema técnico? ¿Qué hicieron para resolverlo?}
  \resp Un archivo con errores dejaba el safari a medio cargar, porque las entidades se registran en el tablero desde su constructor. Se resolvió guardando una referencia al tablero anterior y restaurándola si la importación falla.

  \item \pregunta{¿Qué hicieron bien como equipo? ¿Qué se comprometen a hacer para mejorar?}
  \resp Se dividió bien el trabajo entre modelo, vista y pruebas. Como mejora, se comprometen a escribir las pruebas de unidad antes de implementar cada método.

  \item \pregunta{¿Qué referencias usaron? ¿Cuál fue la más útil?}
  \resp La más útil fue la documentación oficial de Java sobre serialización y reflexión.
\end{enumerate}

% ------------------------------------------------------------------
\begin{thebibliography}{9}
\setlength{\itemsep}{0.3em}

\bibitem{serial}
Oracle. (2024). \textit{Java Object Serialization Specification}. Oracle Corporation. \url{https://docs.oracle.com/en/java/javase/21/docs/specs/serialization/index.html}

\bibitem{reflect}
Oracle. (2024). \textit{Class Class<T>} [Documentación de la API de Java SE 21]. \url{https://docs.oracle.com/en/java/javase/21/docs/api/java.base/java/lang/Class.html}

\bibitem{io}
Oracle. (2024). \textit{The Java Tutorials: Basic I/O}. \url{https://docs.oracle.com/javase/tutorial/essential/io/}

\bibitem{swing}
Oracle. (2024). \textit{How to Use File Choosers}. The Java Tutorials. \url{https://docs.oracle.com/javase/tutorial/uiswing/components/filechooser.html}

\bibitem{xp}
Beck, K., \& Andres, C. (2004). \textit{Extreme Programming Explained: Embrace Change} (2.ª ed.). Addison-Wesley.

\end{thebibliography}

\end{document}
